\documentclass[11pt]{article}

\usepackage{epsfig}
\usepackage{amsfonts}
\usepackage{amssymb}
\usepackage{amstext}
\usepackage{amsmath}
\usepackage{xspace}
\usepackage{theorem}
\usepackage{times}
\usepackage{graphicx}

%\usepackage{layout}% if you want to see the layout parameters
                     % and now use \layout command in the body

% This is the stuff for normal spacing
\makeatletter
 \setlength{\textwidth}{6.5in}
 \setlength{\oddsidemargin}{0in}
 \setlength{\evensidemargin}{0in}
 \setlength{\topmargin}{0.25in}
 \setlength{\textheight}{8.25in}
 \setlength{\headheight}{0pt}
 \setlength{\headsep}{0pt}
 \setlength{\marginparwidth}{59pt}

 \setlength{\parindent}{0pt}
 \setlength{\parskip}{5pt plus 1pt}
 \setlength{\theorempreskipamount}{5pt plus 1pt}
 \setlength{\theorempostskipamount}{0pt}
 \setlength{\abovedisplayskip}{8pt plus 3pt minus 6pt}

 \renewcommand{\section}{\@startsection{section}{1}{0mm}%
                                   {2ex plus -1ex minus -.2ex}%
                                   {1.3ex plus .2ex}%
                                   {\normalfont\Large\bfseries}}%
 \renewcommand{\subsection}{\@startsection{subsection}{2}{0mm}%
                                     {1ex plus -1ex minus -.2ex}%
                                     {1ex plus .2ex}%
                                     {\normalfont\large\bfseries}}%
 \renewcommand{\subsubsection}{\@startsection{subsubsection}{3}{0mm}%
                                     {1ex plus -1ex minus -.2ex}%
                                     {1ex plus .2ex}%
                                     {\normalfont\normalsize\bfseries}}
 \renewcommand{\paragraph}{\@startsection{paragraph}{4}{0mm}%
                                    {1ex \@plus1ex \@minus.2ex}%
                                    {-1em}%
                                    {\normalfont\normalsize\bfseries}}
 \renewcommand{\subparagraph}{\@startsection{subparagraph}{5}{\parindent}%
                                       {2.0ex \@plus1ex \@minus .2ex}%
                                       {-1em}%
                                      {\normalfont\normalsize\bfseries}}
\makeatother

\newenvironment{proof}{{\bf Proof:  }}{\hfill\rule{2mm}{2mm}}
\newenvironment{proofof}[1]{{\bf Proof of #1:  }}{\hfill\rule{2mm}{2mm}}
\newenvironment{proofofnobox}[1]{{\bf#1:  }}{}
\newenvironment{example}{{\bf Example:  }}{\hfill\rule{2mm}{2mm}}
\renewcommand{\thesection}{\arabic{section}}

\renewcommand{\theequation}{\thesection.\arabic{equation}}
\renewcommand{\thefigure}{\thesection.\arabic{figure}}

\newtheorem{fact}{Fact}[section]
\newtheorem{lemma}[fact]{Lemma}
\newtheorem{theorem}[fact]{Theorem}
\newtheorem{definition}[fact]{Definition}
\newtheorem{corollary}[fact]{Corollary}
\newtheorem{proposition}[fact]{Proposition}
\newtheorem{claim}[fact]{Claim}
\newtheorem{exercise}[fact]{Exercise}

% math notations
\newcommand{\R}{\ensuremath{\mathbb R}}
\newcommand{\Z}{\ensuremath{\mathbb Z}}
\newcommand{\N}{\ensuremath{\mathbb N}}
\newcommand{\F}{\ensuremath{\mathcal F}}
\newcommand{\SymGrp}{\ensuremath{\mathfrak S}}

\newcommand{\size}[1]{\ensuremath{\left|#1\right|}}
\newcommand{\ceil}[1]{\ensuremath{\left\lceil#1\right\rceil}}
\newcommand{\floor}[1]{\ensuremath{\left\lfloor#1\right\rfloor}}
\newcommand{\poly}{\operatorname{poly}}
\newcommand{\polylog}{\operatorname{polylog}}

% asymptotic notations
\newcommand{\Oh}[1]{{\mathcal O}\left({#1}\right)}
\newcommand{\LOh}[1]{{\mathcal O}\left({#1}\right.}
\newcommand{\ROh}[1]{{\mathcal O}\left.{#1}\right)}
\newcommand{\oh}[1]{{o}\left({#1}\right)}
\newcommand{\Om}[1]{{\Omega}\left({#1}\right)}
\newcommand{\om}[1]{{\omega}\left({#1}\right)}
\newcommand{\Th}[1]{{\Theta}\left({#1}\right)}


% pseudocode notations
\newcommand{\xif}{{\bf{\em{if~}}}}
\newcommand{\xthen}{{\bf{\em{then~}}}}
\newcommand{\xelse}{{\bf{\em{else~}}}}
\newcommand{\xelseif}{{\bf{\em{elif~}}}}
\newcommand{\xfi}{{\bf{\em{fi~}}}}
\newcommand{\xcase}{{\bf{\em{case~}}}}
\newcommand{\xendcase}{{\bf{\em{endcase~}}}}
\newcommand{\xfor}{{\bf{\em{for~}}}}
\newcommand{\xto}{{\bf{\em{to~}}}}
\newcommand{\xby}{{\bf{\em{by~}}}}
\newcommand{\xdownto}{{\bf{\em{downto~}}}}
\newcommand{\xdo}{{\bf{\em{do~}}}}
\newcommand{\xrof}{{\bf{\em{rof~}}}}
\newcommand{\xwhile}{{\bf{\em{while~}}}}
\newcommand{\xendwhile}{{\bf{\em{endwhile~}}}}
\newcommand{\xand}{{\bf{\em{and~}}}}
\newcommand{\xor}{{\bf{\em{or~}}}}
\newcommand{\xerror}{{\bf{\em{error~}}}}
\newcommand{\xreturn}{{\bf{\em{return~}}}}
\newcommand{\xparallel}{{\bf{\em{parallel~}}}}
\newcommand{\T}{\hspace{0.5cm}}
\newcommand{\m}{\mathcal}

\def\sland{~\land~}
\def\slor{~\lor~}
\def\sRightarrow{~\Rightarrow~}

\def\comment#1{\hfill{$\left\{\textrm{{\em{#1}}}\right\}$}}
\def\lcomment#1{\hfill{$\left\{\textrm{{\em{#1}}}\right.$}}
\def\rcomment#1{\hfill{$\left.\textrm{{\em{#1}}}\right\}$}}
\def\fcomment#1{\hfill{$\textrm{{\em{#1}}}$}}
\def\func#1{\textrm{\bf{\sc{#1}}}}
\def\funcbf#1{\textrm{\textbf{\textsc{#1}}}}

\newcommand{\hide}[1]{}

\newcommand{\prob}[1]{\ensuremath{\text{{\bf Pr}$\left[#1\right]$}}}
\newcommand{\expct}[1]{\ensuremath{\text{{\bf E}$\left[#1\right]$}}}
\newcommand{\Event}{{\mathcal E}}

\newcommand{\mnote}[1]{\normalmarginpar \marginpar{\tiny #1}}

\makeatletter
   \newcommand\figcaption{\def\@captype{figure}\caption}
   \newcommand\tabcaption{\def\@captype{table}\caption}
\makeatother


%%%%%%%%%%%%%%%%%%%%%%%%%%%%%%%%%%%%%%%%%%%%%%%%%%%%%%%%%%%%%%%%%%%%%%%%%%%
% Document begins here %%%%%%%%%%%%%%%%%%%%%%%%%%%%%%%%%%%%%%%%%%%%%%%%%%%%
%%%%%%%%%%%%%%%%%%%%%%%%%%%%%%%%%%%%%%%%%%%%%%%%%%%%%%%%%%%%%%%%%%%%%%%%%%%

\newcommand{\headings}[4]{
{\bf CSE548: Analysis of Algorithms, Fall 2026} \hfill {{\bf Lecturer:} #1}\\
{{\bf Topic:} #2} \hfill {{\bf Date:} #3} \\
{{\bf Scribe:} #4}\\
\rule[0.1in]{\textwidth}{0.025in}
%\thispagestyle{empty}
}

\begin{document}

\headings{Rezaul A. Chowdhury}{Generating Functions}{October 21, 2026}{Lichen Yang}


The lecture continued the study of generating functions: it revisited a constrained
counting problem, derived a closed form for Fibonacci numbers, and began the
average-case analysis of deterministic Quicksort. The main technique is to encode
a sequence as a power series, manipulate that series algebraically, and recover
information about the sequence from its coefficients.

\section{Generating Functions and a Counting Problem}

\subsection{Representing a sequence}

The ordinary generating function of a sequence $s_0,s_1,s_2,\ldots$ is the formal
power series
\[
 S(z)=\sum_{n=0}^{\infty}s_nz^n
     =s_0+s_1z+s_2z^2+\cdots.
\]
Thus $s_n$ is the coefficient of $z^n$, written $[z^n]S(z)$. The variable $z$
keeps track of the index; in a counting problem it records the number of objects
selected. We manipulate these series formally, so convergence is not needed for
coefficient identities. If interpreted as a numerical series, the geometric
identity below holds for $|z|<1$:
\[
 \frac{1}{1-z}=\sum_{n=0}^{\infty}z^n.
\]
Differentiating term by term and shifting the index gives the useful identity
\begin{equation}
 \frac{1}{(1-z)^2}
 =\sum_{n=1}^{\infty}nz^{n-1}
 =\sum_{n=0}^{\infty}(n+1)z^n.
 \label{eq:geometric-derivative}
\end{equation}

\subsection{The fruit-selection problem}

How many ways are there to buy exactly $n$ fruits under the following constraints?
The two apples are distinguishable by color; for the other fruits, only the
quantity selected matters.
\begin{itemize}
 \item \textbf{Apples:} one red and one green apple are available. There is one
 way to take none, two ways to take one, and one way to take both:
 \[
  A(z)=1+2z+z^2=(1+z)^2.
 \]
 \item \textbf{Bananas:} only three bananas are not rotten, so one may take
 $0,1,2$, or $3$:
 \[
  B(z)=1+z+z^2+z^3=\frac{1-z^4}{1-z}.
 \]
 \item \textbf{Figs:} any number of packs of six may be purchased:
 \[
  F(z)=1+z^6+z^{12}+\cdots=\frac{1}{1-z^6}.
 \]
 \item \textbf{Mangoes:} they are sold in pairs, with at most two pairs allowed:
 \[
  M(z)=1+z^2+z^4=\frac{1-z^6}{1-z^2}.
 \]
 \item \textbf{Peaches:} any number of packs of four may be purchased:
 \[
  P(z)=1+z^4+z^8+\cdots=\frac{1}{1-z^4}.
 \]
\end{itemize}

Let $s_n$ count the valid selections of $n$ fruits. Independent choices are
combined by multiplying their generating functions. For example, a choice of
$a$ apples and $b$ bananas contributes to the coefficient of $z^{a+b}$, and
the number of such combined choices is the product of the two coefficients.
Summing over all possible quantities is exactly what polynomial multiplication
does. Therefore
\begin{align*}
 S(z)&=A(z)B(z)F(z)M(z)P(z)\\
 &= (1+z)^2\frac{1-z^4}{1-z}\frac{1}{1-z^6}
       \frac{1-z^6}{1-z^2}\frac{1}{1-z^4}\\
 &=\frac{(1+z)^2}{(1-z)(1-z^2)}
  =\frac{1+z}{(1-z)^2}.
\end{align*}
Using equation~\eqref{eq:geometric-derivative},
\[
 S(z)=(1+z)\sum_{n=0}^{\infty}(n+1)z^n
     =\sum_{n=0}^{\infty}(2n+1)z^n.
\]
For $n\geq 1$, the two contributions to the coefficient are $n+1$ and $n$;
for $n=0$, the coefficient is $1$. Hence in every case
\[
 \boxed{s_n=2n+1\qquad(n\geq 0).}
\]

\section{Fibonacci Numbers}

\subsection{Encoding the recurrence}

The Fibonacci sequence is defined by
\[
 f_0=0,\qquad f_1=1,\qquad f_n=f_{n-1}+f_{n-2}\quad(n\geq 2),
\]
so its first terms are $0,1,1,2,3,5,8,\ldots$. For convenience set $f_n=0$
for $n<0$. With \emph{Iverson bracket} notation,
\[
 [E]=\begin{cases}1,&\text{if the statement $E$ is true},\\
                  0,&\text{otherwise},\end{cases}
\]
all initial conditions and the recurrence can be written as a single formula:
\begin{equation}
 f_n=f_{n-1}+f_{n-2}+[n=1]\qquad(n\geq 0).
 \label{eq:fib-recurrence}
\end{equation}
For $n=0$ the right side is zero; for $n=1$ the bracket supplies the initial
value $1$; for $n\geq2$ the usual recurrence remains.

In this section, let $F(z)$ denote the generating function of the Fibonacci
sequence (rather than the fig-selection function used above):
\[
 F(z)=\sum_{n=0}^{\infty}f_nz^n.
\]
Multiply equation~\eqref{eq:fib-recurrence} by $z^n$ and sum over $n\geq 0$:
\begin{align*}
 F(z)
 &=\sum_{n=0}^{\infty}f_{n-1}z^n
   +\sum_{n=0}^{\infty}f_{n-2}z^n
   +\sum_{n=0}^{\infty}[n=1]z^n\\
 &=\sum_{n=1}^{\infty}f_{n-1}z^n
   +\sum_{n=2}^{\infty}f_{n-2}z^n+z\\
 &=z\sum_{m=0}^{\infty}f_mz^m
   +z^2\sum_{m=0}^{\infty}f_mz^m+z\\
 &=zF(z)+z^2F(z)+z.
\end{align*}
The omitted terms vanish because negative-index Fibonacci numbers are zero.
Rearranging yields
\begin{equation}
 \boxed{F(z)=\frac{z}{1-z-z^2}.}
 \label{eq:fib-gf}
\end{equation}
This is a closed form for the generating function. To obtain a closed form for
$f_n$, we still need to extract its coefficient of $z^n$.

\subsection{Factoring and extracting coefficients}

Write
\[
 1-z-z^2=(1-\varphi z)(1-\widehat\varphi z).
\]
Comparing coefficients gives
\[
 \varphi+\widehat\varphi=1,\qquad
 \varphi\widehat\varphi=-1.
\]
Substituting $\widehat\varphi=1-\varphi$ gives
$\varphi^2-\varphi-1=0$. By the quadratic formula, the two values are
\[
 \varphi=\frac{1+\sqrt5}{2},\qquad
 \widehat\varphi=\frac{1-\sqrt5}{2},\qquad
 \varphi-\widehat\varphi=\sqrt5.
\]
Here $\varphi\approx1.618$ is the golden ratio and
$\widehat\varphi\approx-0.618$. Partial fractions turn
 equation~\eqref{eq:fib-gf} into geometric series:
\begin{align*}
 F(z)
 &=\frac{1}{\varphi-\widehat\varphi}
   \left(\frac{1}{1-\varphi z}
             -\frac{1}{1-\widehat\varphi z}\right)\\
 &=\frac{1}{\sqrt5}
   \left(\sum_{n=0}^{\infty}\varphi^n z^n
             -\sum_{n=0}^{\infty}\widehat\varphi^{\,n}z^n\right)\\
 &=\sum_{n=0}^{\infty}
       \frac{\varphi^n-\widehat\varphi^{\,n}}{\sqrt5}z^n.
\end{align*}
Equating coefficients proves the following formula.
\begin{theorem}[Closed form for Fibonacci numbers]
For every integer $n\geq0$,
\[
 \boxed{f_n=\frac{\varphi^n-\widehat\varphi^{\,n}}{\sqrt5}
 =\frac{1}{\sqrt5}\left(\frac{1+\sqrt5}{2}\right)^n
  -\frac{1}{\sqrt5}\left(\frac{1-\sqrt5}{2}\right)^n.}
\]
\end{theorem}

\subsection{Why the iterative algorithm takes quadratic bit time}

The straightforward algorithm computes the sequence in increasing index order:
\begin{quote}
\func{Fibonacci$(n)$}\\
\xif $n=0$ \xthen \xreturn $0$\\
$f_0\leftarrow0$; $f_1\leftarrow1$\\
\xfor $i\leftarrow2$ \xto $n$ \xdo\\
\T $f_i\leftarrow f_{i-1}+f_{i-2}$\\
\xreturn $f_n$
\end{quote}
There are $\Theta(n)$ additions, but they are not constant-time operations on
arbitrarily large integers. Since $|\widehat\varphi|<1$ and $\varphi>1$,
\[
 f_n\sim\frac{\varphi^n}{\sqrt5},\qquad
 \log_2 f_n=n\log_2\varphi-\log_2\sqrt5+o(1)=\Theta(n).
\]
Consequently, the binary representation of $f_n$ has $\Theta(n)$ bits.
At iteration $i$, ordinary binary addition operates on $\Theta(i)$-bit integers
and takes $\Theta(i)$ bit operations. Thus the total cost is
\[
 \sum_{i=2}^{n}\Theta(i)=\boxed{\Theta(n^2)}.
\]
Only the two preceding values need to be retained, so storing the entire array
is unnecessary. This saves storage but does not change the number or size of
the additions.

The closed form likewise does not give a constant-time exact algorithm:
merely writing the binary answer requires $\Omega(n)$ time. Repeated squaring
can reduce the number of multiplications needed for exponentiation to
$O(\log n)$, but the sizes of the operands and the precision needed for an exact
answer still matter. The lecture left faster computation as a question to
consider and noted that Fibonacci numbers would reappear in the analysis of
Fibonacci heaps.

\section{Average-Case Analysis of Deterministic Quicksort}

\subsection{The algorithm and stable partitioning}

The input is an array $A[1:n]$ of $n$ distinct, comparable values. The goal is
to rearrange them in increasing order. The variant considered in this lecture
always uses the \emph{first} element as its pivot.
\begin{enumerate}
 \item If $n\leq1$, return: the array is already sorted.
 \item Choose the pivot $x=A[1]$.
 \item Compare each of the other $n-1$ elements with $x$ and stably partition
 the array as $L,x,R$, where all values in $L$ are smaller than $x$ and all
 values in $R$ are larger.
 \item Recursively sort $L$ and $R$.
\end{enumerate}
If the pivot has rank $k$, it ends at position $k$, and the two recursive
subproblems have sizes $k-1$ and $n-k$. After both are sorted, the entire array
is sorted; no merging step is needed.

\textbf{Stable partitioning} means that two elements which end up in the same
partition retain their original relative order. Elements ending up on opposite
sides need not retain that order. A stable partition is straightforward with
an auxiliary array. The analysis below counts element comparisons; it does not
assume that this particular stable implementation uses constant auxiliary space.

\subsection{Best and worst cases}

The split sizes determine the shape of the recursion tree. If every pivot is
the smallest or largest remaining element, one side is empty and the other has
one fewer element. For example, an increasing input has this behavior when the
first element is always chosen. Its comparison count is
\[
 (n-1)+(n-2)+\cdots+1=\frac{n(n-1)}{2}=\Theta(n^2).
\]
When the splits are balanced, the tree has $\Theta(\log n)$ levels, with at
most $O(n)$ partitioning work per level, and the total is $\Theta(n\log n)$.
The corresponding balanced recurrence is
\[
 B(n)=B\!\left(\left\lfloor\frac{n-1}{2}\right\rfloor\right)
     +B\!\left(\left\lceil\frac{n-1}{2}\right\rceil\right)+(n-1).
\]
Thus Quicksort has quadratic worst-case behavior but $\Theta(n\log n)$
best-case behavior. Its good practical performance motivates asking what
happens \emph{on average}. The lecturer stated that the average is also
$\Theta(n\log n)$ and began a generating-function derivation of this result.

\subsection{Reducing the inputs to permutations of ranks}

Quicksort is comparison-based. Replacing every input value by its rank
preserves every comparison result, hence every branch and the number of
comparisons performed. For example,
\[
 (40,10,30,20)\quad\longmapsto\quad(4,1,3,2)
\]
produces the same comparison behavior. This replacement is a device for the
analysis, not a preprocessing step performed by the algorithm.

We may therefore restrict attention to the $n!$ permutations of
$\{1,\ldots,n\}$. Let $C(\pi)$ denote the number of comparisons on permutation
$\pi$, and let $\mathfrak S_n$ be the set of these permutations. Define
\[
 t_n=\frac{1}{n!}\sum_{\pi\in\mathfrak S_n}C(\pi),\qquad t_0=t_1=0.
\]
This is an arithmetic average over inputs for a \emph{deterministic} algorithm.
There is no random pivot choice. In particular, the averaging factor derived
below comes from counting permutations, not from random choices made during
execution. The analysis of randomized Quicksort was deferred to a later lecture.

\subsection{Deriving the recurrence}

Fix the first element to be $k$. The remaining $n-1$ values can appear in
$(n-1)!$ orders, so exactly $(n-1)!$ input permutations use $k$ as the pivot.
Partitioning costs $n-1$ comparisons and leaves subproblems of sizes $k-1$
and $n-k$.

To justify using $t_{k-1}$ and $t_{n-k}$ for the recursive costs, consider any
fixed ordering of the $k-1$ smaller elements and any fixed ordering of the
$n-k$ larger elements. There are
\[
 \binom{n-1}{k-1}
\]
ways to interleave those two orders after the pivot. Stable partitioning
recovers those same two orders. This number is independent of the chosen
orders, so each ordering of either subproblem occurs equally often when we
average over the inputs beginning with $k$. Their average recursive costs are
therefore $t_{k-1}$ and $t_{n-k}$.

Summing the costs for all possible first elements and dividing by $n!$ gives,
for $n\geq1$,
\begin{align*}
 t_n
 &=\frac{1}{n!}\sum_{k=1}^{n}(n-1)!
       \bigl((n-1)+t_{k-1}+t_{n-k}\bigr)\\
 &=\frac{1}{n}\sum_{k=1}^{n}
       \bigl((n-1)+t_{k-1}+t_{n-k}\bigr)\\
 &=n-1+\frac{1}{n}\sum_{k=1}^{n}(t_{k-1}+t_{n-k}).
\end{align*}
Here $(n-1)!/n!=1/n$. Also, $\sum_{k=1}^{n}(n-1)=n(n-1)$:
the summand $n-1$ is constant with respect to $k$, so this is not the
sum $1+2+\cdots+(n-1)$.

As $k$ goes from $1$ to $n$, the terms $t_{k-1}$ run from $t_0$ to $t_{n-1}$,
while $t_{n-k}$ runs through exactly the same list in reverse. Thus
\begin{equation}
 \boxed{t_0=0,\qquad
 t_n=n-1+\frac{2}{n}\sum_{k=0}^{n-1}t_k\quad(n\geq1).}
 \label{eq:quicksort-recurrence}
\end{equation}
The recurrence itself gives $t_1=0$. For example, $t_2=1$ and
$t_3=2+\frac23(t_0+t_1+t_2)=\frac83$, consistent with averaging over the
six permutations of three values.

\subsection{Solving the recurrence with generating functions}

Let
\[
 T(z)=\sum_{n=0}^{\infty}t_nz^n.
\]
Since $t_0=0$, substituting equation~\eqref{eq:quicksort-recurrence} gives
\[
 T(z)=t_0+\sum_{n=1}^{\infty}
       \left(n-1+\frac{2}{n}\sum_{k=0}^{n-1}t_k\right)z^n.
\]
The factor $1/n$ suggests differentiating: the derivative of $z^n$ supplies
exactly the factor $n$ needed to cancel it. This yields
\begin{equation}
 T'(z)=\sum_{n=1}^{\infty}
       \left(n(n-1)+2\sum_{k=0}^{n-1}t_k\right)z^{n-1}.
 \label{eq:quicksort-derivative}
\end{equation}
The index $n$ varies inside the sum, so one cannot simply multiply the whole
generating-function identity by a single $n$. Differentiation performs the
required multiplication separately on each coefficient.

To complete the analysis, the geometric-series identities simplify
 equation~\eqref{eq:quicksort-derivative} to
\[
 T'(z)=\frac{2z}{(1-z)^3}+\frac{2T(z)}{1-z},
 \qquad
 \frac{d}{dz}\bigl((1-z)^2T(z)\bigr)=\frac{2z}{1-z}.
\]
Here the prefix sums contribute $T(z)/(1-z)$, since
$\sum_{n\geq1}\sum_{k=0}^{n-1}t_kz^{n-1}
=\sum_{k\geq0}t_kz^k\sum_{j\geq0}z^j$.
Integrating and using $T(0)=t_0=0$ gives
\[
 T(z)=\frac{-2\ln(1-z)-2z}{(1-z)^2}
     =2\left(\sum_{k=2}^{\infty}\frac{z^k}{k}\right)
        \left(\sum_{j=0}^{\infty}(j+1)z^j\right).
\]
Extracting the coefficient of $z^n$, for $n\geq1$, yields
\[
 t_n=2\sum_{k=2}^{n}\frac{n+1-k}{k}
     =2(n+1)H_n-4n,
 \qquad H_n=\sum_{k=1}^{n}\frac1k.
\]
The harmonic-number bounds $\ln(n+1)\leq H_n\leq1+\ln n$ imply
$H_n=\ln n+O(1)$. Therefore
\[
 \boxed{t_n=2(n+1)H_n-4n=2n\ln n+O(n)=\Theta(n\log n).}
\]
Thus deterministic Quicksort performs $\Theta(n\log n)$ comparisons on average
 over all input permutations, despite its $\Theta(n^2)$ worst case.

\end{document}
